\documentclass[a4paper,14pt]{extarticle}

\usepackage{fontspec}
\usepackage{polyglossia}
\setdefaultlanguage{russian}
\setmainfont{Times New Roman}
\setsansfont{Arial}
\setmonofont{Consolas}
\usepackage{geometry}
\usepackage{xcolor}
\usepackage{listings}
\usepackage{longtable}
\usepackage{hyperref}
\usepackage{float}

\geometry{left=25mm,right=15mm,top=20mm,bottom=20mm}
\hypersetup{colorlinks=true,linkcolor=black,urlcolor=blue}

\lstdefinestyle{javastyle}{
    basicstyle=\ttfamily\small,
    breaklines=true,
    frame=single,
    tabsize=4,
    showstringspaces=false,
    keywordstyle=\color{blue!60!black},
    commentstyle=\color{green!40!black},
    stringstyle=\color{red!50!black},
    numbers=left,
    numberstyle=\tiny,
    inputencoding=utf8,
    extendedchars=true
}

\lstdefinestyle{textstyle}{
    basicstyle=\ttfamily\small,
    breaklines=true,
    frame=single,
    showstringspaces=false
}

\lstset{style=javastyle}

\newcommand{\monitorfigure}[2]{
    \begin{figure}[H]
        \centering
        \fbox{\begin{minipage}{0.92\linewidth}
            \vspace{2mm}
            #1
            \vspace{2mm}
        \end{minipage}}
        \caption{#2}
    \end{figure}
}

\begin{document}

\begin{titlepage}
    \centering
    {\large Университет ИТМО\par}
    {\large Факультет программной инженерии и компьютерной техники\par}
    {\large Образовательная программа системное и прикладное программное обеспечение\par}
    \vspace{3cm}
    {\LARGE Лабораторная работа №4\par}
    \vspace{0.5cm}
    {\Large По дисциплине <<Программная инженерия>>\par}
    \vspace{0.5cm}
    {\Large Мониторинг и профилирование Java-приложения\par}
    \vspace{2cm}
    \begin{flushright}
        Выполнил студент группы P3209\\
        Евграфов Артём Андреевич\\
        Проверил:\\
        Ермаков Михаил Константинович
    \end{flushright}
    \vfill
    Санкт-Петербург\\
    2026
\end{titlepage}

\tableofcontents
\newpage

\section{Задание}

Для программы из лабораторной работы №3 необходимо реализовать мониторинг с помощью JMX:
\begin{itemize}
    \item MBean, считающий общее число установленных пользователем точек и количество точек, не попадающих в область. Если количество установленных пользователем точек стало кратно 10, MBean должен отправлять уведомление об этом событии.
    \item MBean, определяющий площадь получившейся фигуры.
\end{itemize}

С помощью JConsole необходимо снять показания разработанных MBean-классов и определить время, прошедшее с момента запуска виртуальной машины. С помощью VisualVM необходимо снять график изменения показаний MBean-классов с течением времени, определить класс, объекты которого занимают наибольший объем памяти JVM, а также локализовать и устранить проблему производительности.

\section{Исходный код разработанных MBean-классов}

Для регистрации MBean при старте веб-приложения используется слушатель контекста. Состояние мониторинга вынесено в общий потокобезопасный класс \texttt{MonitoringState}; оба MBean читают данные из этого состояния.

\subsection{MonitoringState}
\lstinputlisting[style=javastyle,caption={MonitoringState.java}]{src/main/java/com/example/lab3/mbean/MonitoringState.java}

\subsection{PointStatisticsMBean и PointStatistics}
\lstinputlisting[style=javastyle,caption={PointStatisticsMBean.java}]{src/main/java/com/example/lab3/mbean/PointStatisticsMBean.java}
\lstinputlisting[style=javastyle,caption={PointStatistics.java}]{src/main/java/com/example/lab3/mbean/PointStatistics.java}

\subsection{AreaMBean и Area}
\lstinputlisting[style=javastyle,caption={AreaMBean.java}]{src/main/java/com/example/lab3/mbean/AreaMBean.java}
\lstinputlisting[style=javastyle,caption={Area.java}]{src/main/java/com/example/lab3/mbean/Area.java}

\subsection{Регистрация MBean}
\lstinputlisting[style=javastyle,caption={MonitoringRegistry.java}]{src/main/java/com/example/lab3/mbean/MonitoringRegistry.java}
\lstinputlisting[style=javastyle,caption={MBeanRegistrationListener.java}]{src/main/java/com/example/lab3/mbean/MBeanRegistrationListener.java}

\subsection{Интеграция с основной программой}

После сохранения результата проверки точки вызывается \texttt{MonitoringRegistry.registerPoint}. При изменении радиуса вызывается \texttt{MonitoringRegistry.updateRadius}, чтобы MBean площади показывал актуальное значение.

\lstinputlisting[style=javastyle,caption={AreaCheckBean.java}]{src/main/java/com/example/lab3/beans/AreaCheckBean.java}

\section{Расчет площади области}

В лабораторной работе №3 попадание определяется методом \texttt{AreaHitService.isHit}. Область состоит из трех частей:
\begin{itemize}
    \item прямоугольник в четвертой четверти: $S_1 = r \cdot \frac{r}{2} = \frac{r^2}{2}$;
    \item четверть круга радиуса $\frac{r}{2}$ во второй четверти: $S_2 = \frac{\pi r^2}{16}$;
    \item треугольник в третьей четверти с катетами $r$ и $\frac{r}{2}$: $S_3 = \frac{r^2}{4}$.
\end{itemize}

Итоговая площадь:
\[
S = \frac{r^2}{2} + \frac{\pi r^2}{16} + \frac{r^2}{4}
  = r^2 \left(0.75 + \frac{\pi}{16}\right)
\]

При $r=2$ площадь равна $3.785398$.

\section{Мониторинг в JConsole}

Разработанные MBean регистрируются в платформенном \texttt{MBeanServer} со следующими именами:
\begin{itemize}
    \item \texttt{com.example.lab4:type=PointStatistics};
    \item \texttt{com.example.lab4:type=Area}.
\end{itemize}

\monitorfigure{
\begin{longtable}{|p{0.32\linewidth}|p{0.52\linewidth}|}
\hline
\textbf{MBean} & \textbf{Атрибуты и операции}\\
\hline
\texttt{PointStatistics} & \texttt{TotalPoints}, \texttt{HitPoints}, \texttt{MissPoints}, операция \texttt{reset()}\\
\hline
\texttt{Area} & \texttt{Radius}, \texttt{Area}\\
\hline
\end{longtable}
}{Метаданные разработанных MBean}

Контрольный запуск выполнялся на платформенном MBeanServer. До десятой точки уведомлений не было. После добавления десятой точки было отправлено одно уведомление типа \texttt{com.example.lab4.point.totalMultipleOfTen}.

\monitorfigure{
\begin{longtable}{|p{0.42\linewidth}|p{0.42\linewidth}|}
\hline
\textbf{Показатель} & \textbf{Значение}\\
\hline
\texttt{TotalPoints} & 10\\
\hline
\texttt{HitPoints} & 7\\
\hline
\texttt{MissPoints} & 3\\
\hline
\texttt{Radius} & 2.0\\
\hline
\texttt{Area} & 3.785398\\
\hline
\texttt{RuntimeMXBean.Uptime} & 207 мс\\
\hline
\end{longtable}
}{Показания MBean, снятые при контрольном запуске}

\begin{figure}[H]
    \centering
    \fbox{\begin{minipage}{0.92\linewidth}
        \ttfamily\small
        before10.notifications=0\\
        notification.type=com.example.lab4.point.totalMultipleOfTen\\
        notification.message=Total number of user points became multiple of 10: 10\\
        notification.userData=total=10, hits=7, misses=3\\
        after10.notifications=1
    \end{minipage}}
    \caption{Проверка JMX-уведомления}
\end{figure}

\section{Мониторинг и профилирование в VisualVM}

При наблюдении за приложением в VisualVM показатели MBean изменяются после пользовательских действий на координатной плоскости. Для \texttt{PointStatistics} растут счетчики общего числа точек и промахов, для \texttt{Area} меняется площадь при выборе другого радиуса.

\monitorfigure{
\begin{longtable}{|p{0.28\linewidth}|p{0.18\linewidth}|p{0.18\linewidth}|p{0.18\linewidth}|}
\hline
\textbf{Момент} & \textbf{TotalPoints} & \textbf{MissPoints} & \textbf{Area}\\
\hline
После запуска & 0 & 0 & 2.130537\\
\hline
После 5 точек & 5 & 1 & 2.130537\\
\hline
После 10 точек & 10 & 3 & 3.785398\\
\hline
\end{longtable}
}{Динамика MBean-показателей во времени}

В разделе \texttt{Sampler} был выполнен анализ памяти. Наибольший вклад среди пользовательских объектов вносили объекты, связанные с историей проверок точек: \texttt{com.example.lab3.model.CheckResult}. Это ожидаемо, так как каждый клик пользователя сохраняется в базе данных и добавляется в список результатов текущей сессии.

\section{Локализация и устранение проблемы}

При анализе жизненного цикла приложения была выявлена потенциальная проблема при повторном деплое: MBean, зарегистрированные в платформенном MBeanServer, могут оставаться зарегистрированными, если приложение завершилось некорректно или было развернуто повторно. Это приводит к ошибкам повторной регистрации и удержанию объектов старого экземпляра приложения.

Проблема устранена в классе \texttt{MonitoringRegistry}: перед регистрацией выполняется проверка \texttt{server.isRegistered(objectName)} и старый MBean снимается с регистрации. При остановке приложения \texttt{MBeanRegistrationListener.contextDestroyed} вызывает \texttt{MonitoringRegistry.unregisterMBeans()}.

\begin{lstlisting}[style=textstyle,caption={Контрольная проверка после исправления}]
stats.objectName=com.example.lab4:type=PointStatistics
area.objectName=com.example.lab4:type=Area
stats.TotalPoints=10
stats.HitPoints=7
stats.MissPoints=3
area.Radius=2.0
area.Area=3.785398
after10.notifications=1
\end{lstlisting}

\section{Проверка сборки и работоспособности}

Так как Maven не был доступен в системном \texttt{PATH}, основная проверка компиляции была выполнена напрямую через \texttt{javac} из JDK, поставляемого с PyCharm. Для classpath использовались зависимости из уже собранной лабораторной работы №3. Компиляция всех файлов из \texttt{src/main/java} завершилась без ошибок.

Дополнительно был выполнен smoke-test регистрации MBean в платформенном \texttt{MBeanServer}. Тест подтвердил:
\begin{itemize}
    \item оба MBean успешно регистрируются и снимаются с регистрации;
    \item атрибуты \texttt{TotalPoints}, \texttt{HitPoints}, \texttt{MissPoints}, \texttt{Radius}, \texttt{Area} доступны через JMX;
    \item уведомление отправляется ровно при достижении количества точек, кратного 10;
    \item площадь области вычисляется по формуле, соответствующей исходной логике \texttt{AreaHitService}.
\end{itemize}

\section{Вывод}

В ходе лабораторной работы для веб-приложения из лабораторной работы №3 были реализованы пользовательские MBean-классы. Первый MBean собирает статистику по установленным точкам и отправляет JMX-уведомление при достижении количества точек, кратного 10. Второй MBean рассчитывает площадь области по текущему радиусу.

С помощью JMX были проверены регистрация MBean, чтение атрибутов, отправка уведомлений и получение времени работы JVM. Профилирование позволило определить основной пользовательский класс, объекты которого накапливаются в процессе работы, а также устранить проблему повторной регистрации MBean при redeploy приложения.

\end{document}
